\chapter{检索增强及智能体}
\label{lab:90_llm_applications}

\noindent\textbf{配套文件：}\path{notebook/90_llm_applications.ipynb}\par
\noindent\textbf{教材关联：}\bookxrefshort{ch:information-retrieval}、\bookxrefshort{ch:agent-architecture}、\bookxrefshort{ch:agent-runtime}、\bookxrefshort{ch:domain-applications}。

\section*{实验目标}
分离检索证据、模型提案与程序执行。

\section*{准备及定位}
先阅读本册实验总则，确认Notebook中的依赖与资产单元。原理计算、库接口迁移和真实模型训练可能具有不同资源要求，应分别记录设备、版本及运行范围。

源码定位可从下列函数开始；应结合前置数据与配置单元阅读。
\begin{itemize}
\item \texttt{my\_tokenize}
\item \texttt{my\_make\_document}
\item \texttt{my\_chunk\_document}
\end{itemize}

\section*{操作及观察}
\begin{enumerate}
\item 固定文档和查询，检查切块、召回、上下文预算及来源身份。
\item 比较检索基线与基于证据的生成，逐条核对回答依据和缺证据处理。
\item 检查工具参数与执行白名单，沿固定工作流和ReAct路径记录状态变化。
\item 对预算耗尽、无效动作及工具失败设置用例，核对终止条件与保留状态。
\end{enumerate}

\section*{结果判读及提交}
提交召回证据、工具轨迹和失败归因。本实验不包含Codex、OpenClaw或Hermes的安装验收；教材介绍的这些平台不等同于此Notebook中的有界执行器。

提交时附上所执行单元范围、环境与制品身份、原始结果及失败说明。将未运行项目明确列出，不能用Notebook中已有输出替代本次执行证据。
